\chapter{Procesos de Markov y simulación}
\label{ch:markov-simulacion}

\section{Dos formas de describir la incertidumbre}
Una cadena de Markov propaga probabilidades entre estados; una simulación
construye realizaciones de eventos y estima medidas por repetición. La primera
puede dar una distribución estacionaria mediante ecuaciones; la segunda admite
redes y reglas más detalladas, a cambio de error de muestreo. QSB conserva
modelos y resultados normalizados distintos para ambos procedimientos
\cite{ross2019introduction,law2015simulation}.

\section{Cadena finita y convención de filas}
\(P_{ij}\) es la probabilidad de pasar del estado \(i\) al \(j\) en un
período. Cada fila suma uno. La distribución es un vector fila:
\[
 s_{t+1}=s_tP,\qquad \pi P=\pi,\qquad \sum_i\pi_i=1.
\]
Se supone transición homogénea en el tiempo y que el próximo estado depende del
actual, sin información adicional del pasado. Los costos \(c_i\) son por
estado y período; su esperanza es \(s_tc\). No son costos de transición.

\subsection{Ejemplo: disponibilidad de una máquina}
\path{manual/examples/prob_markov.json} define Operativa y Parada, costos
\((0,100)\), distribución inicial \((1,0)\), horizonte tres y
\[
 P=\begin{pmatrix}0.9&0.1\\0.4&0.6\end{pmatrix}.
\]
Desde la raíz del repositorio:
\begin{lstlisting}
.build/debug/qsb validate-markov-json manual/examples/prob_markov.json
.build/debug/qsb solve-markov-json manual/examples/prob_markov.json
\end{lstlisting}
Los resultados son:
\begin{center}
\begin{tabular}{rrrr}
\toprule
Período & Operativa & Parada & Costo esperado \\
\midrule
0 & 1 & 0 & 0 \\
1 & 0.9 & 0.1 & 10 \\
2 & 0.85 & 0.15 & 15 \\
3 & 0.825 & 0.175 & 17.5 \\
\bottomrule
\end{tabular}
\end{center}
La distribución estacionaria es \((0.8,0.2)\), con costo esperado \(20\)
por período. Es una media de largo plazo, no un costo acumulado de tres períodos.
Esta cadena es irreducible y aperiódica, por lo que las distribuciones transitorias
convergen a ella. En general, unicidad estacionaria y convergencia desde cualquier
inicio no son equivalentes: una cadena periódica puede oscilar.

QSB resuelve un sistema lineal estacionario y propaga matrices; declara
\texttt{exact} como carácter algorítmico, con aritmética de coma flotante.
El método exige una distribución estacionaria única y devuelve error si su
sistema es singular. No elige una mezcla arbitraria en cadenas con varias
clases cerradas. Si falta distribución inicial, no inventa una: la salida
transitoria está vacía. Validar filas y dimensiones no garantiza que el sistema
estacionario sea resoluble.

En macOS, el editor permite nombres, costos, probabilidades iniciales opcionales,
matriz de transición y horizonte. Verifique la orientación de filas antes de
resolver; el resultado muestra estacionario, propagación, costos e hipótesis.
La edición de estados mantiene dimensiones coherentes; el texto inválido
permanece visible hasta corregirse. JSON sigue disponible para intercambio.

\section{Simulación de eventos discretos}
Una red contiene fuentes de entidades, colas, servidores y compuertas. Las rutas
incluyen probabilidades y tiempos de transferencia; las reglas de servicio pueden
depender del tipo de entidad. QSB admite distribuciones constante, exponencial,
normal, uniforme y triangular. Las muestras normales negativas se recortan
a cero; no se vuelve a muestrear. Ese recorte modifica la media y la
variabilidad respecto de una normal sin recorte. Sus parámetros representan tiempos o valores de
la distribución, no tasas por unidad temporal. Por ejemplo, una media
exponencial \(0.25\) horas equivale a cuatro llegadas/hora.

El backend educativo modela colas FIFO finitas, servidores alternativos,
enrutamiento probabilístico, transferencias y sincronización de ensamblaje.
Sus salidas son estimaciones \texttt{approximate}, aunque una semilla fija
las haga reproducibles. No hay borrado de calentamiento, traza interactiva de
eventos ni garantía de estado estacionario. Tampoco hay backend externo de
simulación disponible.

\subsection{Ejemplo: una ventanilla simulada}
\path{prob_simulacion.json} conecta fuente Cliente, cola Espera de veinte
plazas y servidor Atención. Cada llegada trae una entidad, el tiempo entre
llegadas es exponencial de media \(0.25\) horas, y el servicio exponencial
medio es \(1/6\) horas. Las transferencias duran cero.
\begin{lstlisting}
.build/debug/qsb validate-simulation-json manual/examples/prob_simulacion.json
.build/debug/qsb solve-simulation-json manual/examples/prob_simulacion.json
.build/debug/qsb solve-simulation-json manual/examples/prob_simulacion.json --replications 3
\end{lstlisting}
La ejecución predeterminada dura \(1000\) horas, semilla \(1\), sin
calentamiento eliminado. Genera \(4005\) entidades, completa \(4004\),
no rechaza ninguna, y estima longitud media de cola \(1.202238\), máximo
\(15\), y utilización \(0.664470\). El cliente sin completar puede estar
en cola o servicio al terminar: no es necesariamente rechazado.

Con tres réplicas independientemente sembradas, la longitud media es
\(1.229240\), intervalo normal aproximado del \(95\%\)
\([1.146668,1.311813]\); la utilización es \(0.659950\), intervalo
\([0.638179,0.681722]\). Los contadores agregados publicados son medias
redondeadas: generan \(4006\) y completan \(4005\), no la suma de tres
corridas. Los máximos, en cambio, conservan el máximo entre réplicas.
La semilla base, horizonte y número de réplicas figuran en la salida.
La CLI expone réplicas, pero no opciones para elegir semilla u horizonte en
estos comandos; esos valores proceden de las opciones predeterminadas del núcleo.

La fórmula del intervalo usa \(\bar x\pm1.96s/\sqrt R\). Con pocas
réplicas, la aproximación normal puede ser pobre; aumentar \(R\) reduce el
error de muestreo, no corrige sesgo de arranque ni una distribución incorrecta.
No interprete el intervalo como garantía de que el modelo describe el mundo.
El M/M/1 analítico del capítulo de colas da \(L_q=4/3\) y \(\rho=2/3\).
La simulación tiene cola finita y horizonte finito: la proximidad es útil para
comprobar unidades, pero no exige igualdad exacta. Tres réplicas tampoco
justifican deducir un fallo porque un intervalo no contenga el valor analítico.

\section{Procedimiento en macOS}
En la definición de simulación, complete Componentes, Rutas y Reglas de servicio.
Introduzca los parámetros de cada distribución por separado. Los nombres deben
coincidir con las referencias de rutas y tipos de entidad. Se pueden intercambiar
tablas completas con hojas de cálculo; el pegado comprueba cada sección antes de
reemplazarla. En el panel Run seleccione entre una y cien réplicas,
valide y ejecute. La interfaz muestra estado de cálculo y evita ejecuciones
duplicadas; cambiar el modelo invalida un resultado pendiente. Esto no detiene
un cálculo ya iniciado en el núcleo.

El panel de resultados muestra entidades generadas/completadas, tasa de
completado, horizonte, semilla, cola y utilización, con intervalos cuando están
disponibles. Consulte tablas y JSON además de gráficos. Conserve red, reglas,
semilla, horizonte y réplicas al comparar escenarios: cambiar cualquiera puede
cambiar la interpretación de la diferencia.

\section{Ejercicios}
\begin{enumerate}
\item Quite la distribución inicial del caso Markov y compruebe que queda solo
análisis estacionario. Cambie después la matriz a identidad y observe el error
del sistema singular.
\item Duplique la media de servicio simulada y compare terminaciones, rechazo y
utilización. Explique por qué capacidad de cola no equivale a capacidad de servicio.
\item Repita el caso con diez réplicas y compare amplitud de intervalos. Separe
variabilidad entre corridas de diferencias entre los supuestos de modelos.
\end{enumerate}
